\chapter{Introduction}
\label{ch:intro}

\section{Background}

A modern fulfilment warehouse is a shared workspace. Automated mobile robots move stock between
racking and packing stations while people walk the same floor, and the robot's routing problem is
therefore not only about finding the shortest path. It is about finding the shortest path that is
also safe to execute, and about doing so without causing harm during the period in which it is
still learning.

That last qualification is easy to overlook and is the subject of this report. A policy learned in
a simulator can be evaluated purely on its final quality, because mistakes made during training
cost nothing. A policy learned on a live warehouse floor cannot: every exploratory action is taken
in the real world, and the cost of learning is a cost the operator actually pays.

Reinforcement learning offers two classical algorithms for this class of problem, and they differ
precisely on this point. Q-learning \cite{watkins1989,watkins1992} learns the value of the optimal
policy regardless of how the agent is behaving. SARSA \cite{rummery1994,sutton1996} learns the
value of the policy the agent is actually following, exploration included. On a benign problem the
two converge to the same answer. On a problem where exploration can be costly they do not, and the
difference is both predictable and practically important.

\section{Problem statement}

The task is navigation on a warehouse floor represented as a grid. A robot starts at a charging
station and must reach a delivery point. Racking blocks movement. A pedestrian walkway crosses the
floor, and entering it is a safety incident.

The shortest route runs along the main aisle, which lies directly alongside the walkway. A second
route one row further from the walkway is slightly longer and carries no such exposure. The
problem is interesting because those two routes differ in exactly the way that distinguishes the
two algorithms: the fast route is optimal for an agent that never makes mistakes, and dangerous
for one that is still exploring.

\section{Research question}

\begin{quote}
\textbf{How do Q-learning and SARSA compare when a warehouse robot must learn the shortest
collision-free route to a delivery point?}
\end{quote}

Four subsidiary questions are answered empirically:

\begin{enumerate}
\item \textbf{Deployed quality.} After training, which algorithm's greedy policy produces the
      shorter route, and how close is it to the true optimum?
\item \textbf{Cost of learning.} Which algorithm accumulates fewer safety incidents while learning,
      and what does training cost in accumulated reward?
\item \textbf{Significance.} Are the observed differences distinguishable from seed-to-seed noise?
\item \textbf{Mechanism.} Whatever is observed, what property of the algorithms and of the task
      accounts for it?
\end{enumerate}

\section{Objectives}

\begin{enumerate}
\item Implement a grid-world warehouse environment from first principles, with no reinforcement
      learning library.
\item Implement Q-learning and SARSA from first principles, sharing all machinery except the
      update target, so that the comparison isolates the on-policy/off-policy distinction.
\item Audit the source dataset column by column and use only what carries signal, documenting the
      exclusions with evidence.
\item Verify the environment and both agents numerically before any experiment is run.
\item Compare the algorithms over thirty paired seeds with a paired non-parametric significance
      test and reported effect sizes.
\item Test the robustness of the result against an external baseline, against independently
      generated floors, and against the main modelling assumption.
\end{enumerate}

\section{Scope and delimitations}

\keypoint{Tabular methods only.} The state space is 400 cells, small enough to enumerate. Nothing
here concerns function approximation or deep reinforcement learning, where off-policy bootstrapping
raises stability issues that do not arise in the tabular case.

\keypoint{A simulated environment built on real layouts.} This point is developed at length in
Chapter~\ref{ch:dataset} because it determines what the study can honestly claim. The obstacle
layouts come from the dataset; the dynamics and the reward structure are written here.

\keypoint{One task family.} All results concern grid navigation with a localised hazard. The
report is explicit about which findings depend on that structure.

\keypoint{No tuning contest.} Both algorithms share one hyperparameter setting. The objective is
a controlled comparison, not the best achievable performance from either.

\section{Contributions}

\begin{enumerate}
\item \textbf{A documented dataset audit} showing that several columns of a published Kaggle
      dataset --- its reward signal, its path lengths, its episode structure --- carry no usable
      information, with the evidence rather than the assertion.
\item \textbf{A verification suite for tabular RL}, including an equivalence test that pins the
      two algorithms to each other under a greedy policy, where they are mathematically identical.
\item \textbf{A properly paired multi-seed comparison}, with significance testing and effect
      sizes, rather than a single run per algorithm.
\item \textbf{An explicit separation of what is data from what is assumed}, together with a
      sensitivity analysis over the assumption that drives the result.
\item \textbf{A negative control}: the demonstration that on a floor without a localised hazard,
      the two algorithms converge on the identical optimum and the much-cited distinction between
      them simply does not appear.
\end{enumerate}

\section{Organisation}

Chapter~\ref{ch:literature} reviews temporal-difference control and the on-policy/off-policy
distinction. Chapter~\ref{ch:theory} develops the mathematics: Markov decision processes, the
Bellman equations, TD(0), the two control algorithms and their convergence conditions.
Chapter~\ref{ch:dataset} audits the dataset and specifies the environment, separating data from
assumption. Chapter~\ref{ch:implementation} presents the implementation.
Chapter~\ref{ch:verification} presents the verification suite. Chapter~\ref{ch:results} reports the
experiment. Chapter~\ref{ch:discussion} interprets it and states the threats to validity.
Chapter~\ref{ch:conclusion} concludes.
